\origpage{47}{GRAPH MINORS}
\noindent vertex from $Z_{k_{l'}}$, because $Z_{k_{l'}}$ separates $\{v\colon\mu(v)<k_{l'}\}$ and $\{v\colon\mu(v)\ge k_{l'}\}$.
We may assume therefore that $k_{l'-1}\le\mu(v)<k_{l'}$ for every vertex $v$ of $B_j$. By
the minimality of $k_{l'}$, there is a vertex $v$ of $B_j$ with $\mu(v)\ge k_{l'}-1$; and so
$\mu(v)=k_{l'}-1$. But then $v\in Z_{k_{l'}-1}$, and again our claim is true. Thus for
$2\le l'<l$ we have $|J_{l'}-J_{l'-1}|\le 2m$. A similar argument shows that
$|J_1|\le 2m$. It follows that
\begin{equation*}
|J_{l-1}|\le 2(l-1)m
\end{equation*}
and hence that $J_{l-1}\ne\{1,..., n\}$, since $l\le s=\lceil n/2m\rceil$. Thus $k_l$ is well defined.
Hence by induction, $k_1,..., k_s$ are well defined.

Choose $j_1$ with $1\le j_1\le n$ so that $\mu(v)<k_1$ for every vertex $v$ of $B_{j_1}$; and
for $2\le l\le s$, choose $j_l$ with $1\le j_l\le n$ so that $k_{l-1}\le\mu(v)<k_l$ for every
vertex $v$ of $B_{j_l}$. Put $C_l=B_{j_l}$ $(1\le l\le s)$. Then $C_1,..., C_s$ are disjoint connected
subgraphs of $G$, and each meets at least $\frac12\theta^2$ of $P_1,..., P_m$: and moreover, for
$1\le l<l'\le s$, all vertices of $C_l$ on $P_i$ occur before all vertices of $C_{l'}$. By
(2.4), $s<\theta^{2\theta-2}\binom{m}{\frac12\theta^2}$, and so $n<2m\theta^{2\theta-2}\binom{m}{\frac12\theta^2}$, as required.
\stmt
(2.7). \textit{Suppose that there are $m$ disjoint paths in $G\in\mathscr{F}_\theta$ from $A_1$ to $A_2$,
where $A_1, A_2\subseteq V(G)$; and that $B_1,..., B_n$ are disjoint connected subgraphs of
$G$, and for $1\le j\le n$, $V(B_j)$ separates $A_1$ and $A_2$. Then either $m<\frac12\theta^2$ or
$n<\theta^{2\theta}$.}
\stmt
\textit{Proof}. If possible, choose a counterexample with $|V(G)|+|E(G)|+m$
minimum. Then evidently we have $m=\frac12\theta^2$.

Let $\{P_1,..., P_m\}$ be a set of $m$ disjoint paths in $G$ from $A_1$ to $A_2$. If some $P_i$
and some $B_j$ have an edge in common, we can contract that edge and
produce a smaller counterexample. Thus no edge belongs to both some $P_i$
and some $B_j$. Moreover, every edge belongs either to some $P_i$ or to some $B_j$,
for otherwise we could delete it.

It follows that
\begin{equation*}
E(P_1)\cup\cdots\cup E(P_m)=E(G)-(E(B_1)\cup\cdots\cup E(B_n)).
\end{equation*}
If $\{P_1',..., P_m'\}$ is another set of $m$ disjoint paths from $A_1$ to $A_2$, we have, by
the same argument,
\begin{equation*}
E(P_1')\cup\cdots\cup E(P_m')=E(G)-(E(B_1)\cup\cdots\cup E(B_n))
\end{equation*}
and hence
\begin{equation*}
E(P_1')\cup\cdots\cup E(P_m')=E(P_1)\cup\cdots\cup E(P_m).
\end{equation*}
It follows that $\{P_1',..., P_m'\}=\{P_1,..., P_m\}$, and so the set $\{P_1,..., P_m\}$ is unique.

We claim that every vertex $v$ of $G$ is in some $P_i$. For if not, and $v$ is an

